\subsection*{NT-A12--A14. A nonabelian Frobenius class and a Hecke angle}

Let \(L/K\) be a finite Galois extension with group \(G\), and let
\(C\subset G\) be one conjugacy class. Choose \(g\in C\) and put
\[
 \begin{aligned}
 H&=\langle g\rangle,\quad h=|H|,\quad F=L^H,\quad d=[F:K],\\
 \delta&=\frac{|C|}{|G|},\quad n_F=[F:\mathbf Q],\quad
                              n_L=[L:\mathbf Q]=h n_F .
 \end{aligned}
                                                               \tag{A12.1}
\]
Let \(\eta\) be a unitary Hecke character of \(K\) which is not a
finite-order character times a unitary norm power. Write
\(\nu=\eta\circ N_{F/K}\), and use the primitive character attached
to every product below. Let \(S\) be a finite set of prime ideals of
\(K\) containing the primes ramified in \(L/K\) and the support of
\(\mathfrak f_\eta\). Extra prescribed omissions may be included.
Set
\[
 \begin{aligned}
 \mathfrak s_F&=\prod_{\mathfrak q\mid\mathfrak p,\,
                                      \mathfrak p\in S}\mathfrak q,\\
 B_\nu&=4+\max_j|\kappa_{\nu,j}|,\\
 Q_C&=D_L
       \bigl(N_F\mathfrak f_\nu\,N_F\mathfrak s_F\bigr)^h
       B_\nu^{n_L},\\
 {\cal L}_C(x)&=\log Q_C+2n_L\log(3x).
 \end{aligned}                                                 \tag{A12.2}
\]
Here \(D_L\) is the absolute discriminant; the gamma parameters
\(\kappa_{\nu,j}\) have precisely convention (A08.1).
All conductors and norms in (A12.2) are over \(F\).

The exact analytic input is RH for
\[
 L_F(s,\chi\nu^m),\qquad
       \chi\in\widehat H,\quad m=0,1,2,\ldots,                 \tag{A12.3}
\]
where class field theory regards \(\chi\) as a finite-order Hecke
character for the cyclic extension \(L/F\). In particular the
\((\chi,m)=(1,0)\) member is \(\zeta_F\).
Conditional input \(M\) supplies this family. No holomorphy or RH
for an irreducible nonabelian Artin function is assumed.

\paragraph{NT-A12: uniform joint discrepancy.}
For an arc \(J\subset\mathbf R/\mathbf Z\), count prime ideals
\(\mathfrak p\notin S\) with \(N_K\mathfrak p\le x\),
Frobenius class \(C\), and \(\eta(\mathfrak p)\in e^{2\pi iJ}\).
Then
\[
 N_{C,J}(x)
   =\delta |J|\operatorname{Li}_2(x)
      +O\bigl(\delta\sqrt x\log(2x){\cal L}_C(x)\bigr),
       \qquad x\ge2.                                      \tag{A12.4}
\]
The constant is absolute, uniformly in all the displayed data,
the arc and its endpoint conventions. The same bound holds when
only absolute-degree-one primes of \(K\) are counted.

\begin{proof}
We first check the possible norm degeneracy. The diagonal idele
map \(j:C_K\longrightarrow C_F\) satisfies
\[
 N_{F/K}(j(c))=c^d,\qquad |j(c)|_F=|c|_K^d.
\]
If \(\nu=\theta|\cdot|_F^{it}\) with \(\theta\) of finite order,
composition with \(j\) gives
\((\eta|\cdot|_K^{-it})^d=\theta\circ j\).
If the latter has order dividing \(r\), the former character
raised to \(dr\) is trivial. This contradicts the hypothesis on
\(\eta\). Thus \(\nu\) also is not finite order times norm.
For \(m>0\), an equality
\(\chi\nu^m=|\cdot|_F^{it}\) would make
\((\nu|\cdot|_F^{-it/m})^m=\chi^{-1}\), giving the same
contradiction. Hence every \(m>0\) member of (A12.3) is entire,
with no pure-norm pole. For \(m=0\) only \(\chi=1\) is principal.
The principal member is treated separately by NT-A08.

Next we retain the sum of conductors instead of replacing each
one by its largest value. Locally, the conductor of a product
divides the product of conductors, and taking a power cannot
increase the conductor. Therefore
\[
 \mathfrak f_{\chi\nu^m}\mid\mathfrak f_\chi\mathfrak f_\nu,
 \qquad
 {\cal C}(\chi\nu^m)
   \le D_F N_F\mathfrak f_\chi N_F\mathfrak f_\nu
                       [B_\nu(m+1)]^{n_F}.               \tag{A12.5}
\]
The gamma-factor inequality follows directly from (A08.1):
a finite-order character changes only the real parities, while
the imaginary and complex angular parameters of \(\nu^m\)
grow linearly in \(m\). It also holds at \(m=0\).
The conductor--discriminant formula for the cyclic extension
\(L/F\) gives the exact identity
\[
       D_L=D_F^h\prod_{\chi\in\widehat H}
                                      N_F\mathfrak f_\chi .
                                                               \tag{A12.6}
\]
This includes all ramified local factors. The precise classical
locator is Neukirch, Chapter VII, (11.9), printed p.~534,
with the supplied conductor/induction proof on pp.~533--534
and the one-dimensional conductor identification (11.10),
p.~535; see S-NT-NEUKIRCH. Summing the logarithms in (A12.5)
thus yields
\[
 \sum_{\chi\in\widehat H}\log{\cal C}(\chi\nu^m)
 \le \log D_L+h\log N_F\mathfrak f_\nu
          +n_L\log B_\nu+n_L\log(m+1).                 \tag{A12.7}
\]

Every character in (A12.3) is unramified outside the primes
dividing \(\mathfrak s_F\). Indeed a local norm sends units
to units, and the characters of \(L/F\) ramify only over the
ramification locus already in \(S\). The prime version of
NT-A08, deleting the additional primes in \(\mathfrak s_F\),
therefore gives
\[
 \begin{split}
 P_{\chi,m}(x)
  &:=
    \sum_{\substack{N_F\mathfrak q\le x\\
                      \mathfrak q\nmid\mathfrak s_F}}
                 \chi(\mathfrak q)\nu(\mathfrak q)^m\\
  &=\mathbf1_{\chi=1,m=0}\operatorname{Li}_2(x)
    +O\left(\sqrt x\{\log{\cal C}(\chi\nu^m)
                       +n_F\log(2x)\}
                   +\log N_F\mathfrak s_F\right).
 \end{split}                                                   \tag{A12.8}
\]
Its constant is absolute. The last error is simply the
number of additional omitted prime ideals, at most
\(\log N_F\mathfrak s_F/\log2\); a conductor exponent is
not required to count these omissions. Formula (A12.8)
also follows directly from the proof of (A08.4).

We now establish the induction identity used at primes:
\[
       \mathbf1_C
       =\delta\sum_{\chi\in\widehat H}
                         \overline{\chi(g)}\,
                            \operatorname{Ind}_H^G\chi .
                                                               \tag{A12.9}
\]
Orthogonality on \(H\) says
\(\sum_\chi\overline{\chi(g)}\chi(y)=h\mathbf1_{\{g\}}(y)\).
By the definition of induced class functions,
\(\operatorname{Ind}_H^G\mathbf1_{\{g\}}\) has value
\(|Z_G(g)|/h\) on \(C\) and zero elsewhere. Since
\(\delta |Z_G(g)|=1\), (A12.9) follows.

For a prime \(\mathfrak p\notin S\), the cosets in \(G/H\)
fixed by a Frobenius element correspond exactly to prime
ideals \(\mathfrak q\) of \(F\) over \(\mathfrak p\) with
relative residue degree one. The trace of
\(\operatorname{Ind}_H^G\chi\) is the sum of the corresponding
\(\chi(\operatorname{Frob}_{\mathfrak q})\).
For these primes,
\(\nu(\mathfrak q)=\eta(\mathfrak p)\), because the ideal
norm is \(\mathfrak p\). Use the same arithmetic Frobenius
convention in this identity and in class field reciprocity.
Multiplying (A12.9) by \(\eta(\mathfrak p)^m\) and summing
therefore gives
\[
 T_m^C(x):=
 \sum_{\substack{N_K\mathfrak p\le x,\ \mathfrak p\notin S\\
                  \operatorname{Frob}_{\mathfrak p}\in C}}
                        \eta(\mathfrak p)^m
 =\delta\sum_{\chi\in\widehat H}
              \overline{\chi(g)}P_{\chi,m}(x)
           +O(\delta n_L\sqrt x).                       \tag{A12.10}
\]
To verify the last error, the extra primes in the right-hand
sum have relative residue degree at least two, hence lie
over primes of \(K\) with norm at most \(\sqrt x\).
There are at most \(d[K:\mathbf Q]\sqrt x=n_F\sqrt x\)
such primes in \(F\). Each of the \(h\) character sums has
terms of modulus one, so their total contribution is at
most \(\delta h n_F\sqrt x\). Primes in \(S\) have been
deleted on both sides; no Frobenius value is assigned to
a ramified prime.

Combining (A12.7)--(A12.10), and using
\(\log Q_C\ge n_L\log4\), proves
\[
 T_m^C(x)=\delta\mathbf1_{m=0}\operatorname{Li}_2(x)
     +O\left(\delta\sqrt x
        \{\log Q_C+n_L\log[(m+1)(2x)]\}\right).
                                                               \tag{A12.11}
\]
All constants remain absolute. Notice that the \(h\)
summands have already been absorbed in the exact
discriminant and in \(n_L\); they are not replaced by
an additional factor \(|G|\).

Apply Erd\H{o}s--Tur\'an, S-NT-ET Theorem III, to the angles
of the primes with Frobenius class \(C\), with cutoff
\(\lceil x\rceil\). The zero-frequency truncation error
is \(O([K:\mathbf Q])\), since there are at most
\([K:\mathbf Q]x\) prime ideals of norm at most \(x\).
It is absorbed because
\(\delta n_L=|C|[K:\mathbf Q]\ge[K:\mathbf Q]\).
For the other terms, the braces in (A12.11) are
\(O({\cal L}_C(x))\) throughout \(1\le m\le\lceil x\rceil\);
their harmonic weights sum to \(O(\log(2x))\).
The \(m=0\) instance of (A12.11) supplies the main term
in (A12.4). Decreasing arcs to an endpoint bounds its
singleton count by the same error; thus all endpoint
conventions are allowed.
Finally the number of absolute-degree-at-least-two primes
of \(K\) with norm at most \(x\) is at most
\([K:\mathbf Q]\sqrt x\le\delta n_L\sqrt x\).
Deleting them proves the degree-one version.
\end{proof}

\paragraph{NT-A13: a shrinking angle in the specified class.}
Fix \(c>0\), \(0<\alpha<1/2\) and
\(u_0=\max\{2,c^{1/\alpha}\}\). Fix an angular center.
Let \(N_{C,c,\alpha}(x)\) count exactly the primes
\(\mathfrak p\notin S\) with \(u_0<N_K\mathfrak p\le x\),
Frobenius class \(C\), and angle in the centered arc
of length \(c(N_K\mathfrak p)^{-\alpha}\). Under (A12.3),
\[
 \begin{split}
 N_{C,c,\alpha}(x)
   ={}&\delta c\int_{u_0}^x\frac{u^{-\alpha}}{\log u}\,du\\
     &+O_{c,\alpha}\left(
       \delta\sqrt x\log(2x){\cal L}_C(x)
       +\delta x^{3/4-\alpha/2}\sqrt{{\cal L}_C(x)}
                     \right),\qquad x>\max\{4,u_0\}.
 \end{split}                                                   \tag{A13.1}
\]
The constant is uniform in the center and all the field and
character parameters. A sufficient relative asymptotic regime is
\[
         {\cal L}_C(x)(\log x)^2=o(x^{1/2-\alpha}).
                                                               \tag{A13.2}
\]
For fixed arithmetic data this holds automatically, and the
main term is asymptotic to
\(\delta c x^{1-\alpha}/[(1-\alpha)\log x]\).
The same assertions hold for absolute-degree-one primes.

\begin{proof}
Partition \((u_0,x]\) geometrically with ratio
\(1+\varepsilon\), \(0<\varepsilon\le1/2\).
The centered arcs are nested, so their endpoint arcs bracket
the moving count on each subinterval. Apply (A12.4) to both
endpoint prime counts and sum. Since
\(\sum_j\sqrt{u_j}\ll\sqrt x/\varepsilon\), the total
discrepancy is at most a constant times \(A/\varepsilon\),
where
\[
 A=\delta\sqrt x\log(2x){\cal L}_C(x),\qquad
 B=\frac{\delta c x^{1-\alpha}}{\log x}.
\]
The two main terms bracket the integral in (A13.1).
Their difference is \(O_{c,\alpha}(B\varepsilon)\):
on one mesh interval the ratio of the arc lengths is at most
\((1+\varepsilon)^\alpha=1+O_\alpha(\varepsilon)\), and
\(\int_{u_0}^x u^{-\alpha}du/\log u
 \ll_{c,\alpha}x^{1-\alpha}/\log x\).
The finite initial range in this integral bound depends only
on \(c,\alpha\). Its possible contribution to the count is
absorbed uniformly using
\([K:\mathbf Q]\le\delta n_L\le\delta\log Q_C/\log4\).

Choose \(\varepsilon=\min\{1/2,\sqrt{A/B}\}\).
Then \(A/\varepsilon+B\varepsilon\ll A+\sqrt{AB}\), and
\[
 \sqrt{AB}
  =\delta\sqrt c\,x^{3/4-\alpha/2}
       \sqrt{\frac{\log(2x)}{\log x}{\cal L}_C(x)}.
\]
This proves (A13.1). After division by the main term,
the two errors are bounded, up to fixed \(c,\alpha\)
constants, by
\[
 R(x)={\cal L}_C(x)(\log x)^2x^{\alpha-1/2},
 \qquad \sqrt{R(x)}.
\]
They tend to zero under (A13.2). Applying the identical
partition to the degree-one version of (A12.4) proves
the final assertion.
\end{proof}

\paragraph{NT-A14: a least prime with both conditions.}
For an arc of length \(0<\mu\le1\), put
\[
       V_C=\log(eQ_C)+n_L\log(e n_L/\mu).
\]
There is an absolute constant \(C_0>0\) such that (A12.3)
implies the existence of an absolute-degree-one prime
\(\mathfrak p\notin S\), with Frobenius class \(C\) and
angle in this arc, satisfying
\[
                  N_K\mathfrak p\le C_0\mu^{-2}V_C^6 .
                                                               \tag{A14.1}
\]
No numerical value of \(C_0\) is claimed.

\begin{proof}
Let \(A_0\) be an absolute error constant in (A12.4) for
degree-one primes. Set \(x=C_0\mu^{-2}V_C^6\), with
\(C_0\ge4\). We have \(V_C\ge n_L\),
\(n_L\log(1/\mu)\le V_C\), and
\(n_L\log V_C\le2V_C\).
For the last inequality write \(V_C=n_L v\), note
\(v\ge\log(e n_L)\), and use
\(\log(n_L v)\le2v\). Consequently
\[
 \begin{split}
 \log x&\le(\log C_0+14)V_C/n_L,\\
 \log(2x)&\le(\log C_0+15)V_C/n_L,\\
 {\cal L}_C(x)&\le(32+2\log C_0)V_C .
 \end{split}
\]
Since \(\operatorname{Li}_2(x)\ge x/(2\log x)\),
the main term \(\delta\mu\operatorname{Li}_2(x)\)
exceeds the error in (A12.4) as soon as
\[
 \sqrt{C_0}>
  2A_0(\log C_0+14)(\log C_0+15)(32+2\log C_0).
\]
Here the common factor \(\delta\) cancels; dropping the
remaining factor \(n_L^{-2}\le1\) makes this condition
only stronger. An absolute \(C_0\) satisfying it exists.
The count is then positive, proving (A14.1).
\end{proof}

\paragraph{Geometric meaning and exact work placement.}
The finite condition is a conjugacy-class Frobenius condition
for the finite \'{e}tale \(G\)-cover over the arithmetic open
\(\operatorname{Spec}\mathcal O_K\setminus S\). Passing to
\(F=L^{\langle g\rangle}\) is the corresponding intermediate
cover; its cyclic character decomposition is the explicit
induction (A12.9). The second condition is a rank-one
compact Hecke torus coordinate, pulled back by the norm.
For normalized algebraic Hecke characters this is a finite
Artin factor coupled to a CM torus representation. For
arbitrary unitary infinity type no algebraic-motive
realization is asserted. Formula (A12.3) gives the exact
rank-one analytic family replacing any unproved general
Artin-holomorphy input.
Full proofs: \texttt{frobenius\_hecke\_angles.tex},
NT-A12--A14, also in \texttt{section.tex}.
Run: \texttt{GRH\_NIGHT/number\_theory}.
Classical premises: NT-A08 (S-NT-IK), S-NT-ET and
S-NT-NEUKIRCH; the simultaneous formulations are proved
here, with no literature-priority claim.
